% Body only.  The repository wrapper supplies the title, author, preamble,
% bibliography and document environment.
\newcommand{\RFLC}[1]{L_N[1/3,#1]}
\newcommand{\RFLT}[1]{L_N[2/3,#1]}
\newcommand{\RFF}{\mathbb F}
\newcommand{\RFZ}{\mathbb Z}
\newcommand{\RFQ}{\mathbb Q}
\newcommand{\RFPr}{\mathop{\rm Pr}}
\newcommand{\RFE}{\mathop{\rm E}}

\begin{abstract}
We give an explicit conditional improvement of the recursive randomized
factoring construction in the pinned Cogentic manuscript
\emph{Unconditional Integer Factorization in Time
$L_N[1/3,O(1)]$}.  Its announced constants are $105$ for expected time and
$40$ for deterministic workspace; its printed induction already yields
$104$ and $38$.  We replace dense linear algebra by a verified sparse
kernel algorithm, retain only the hidden roots belonging to accepted
relations, and enlarge the root-prime scale while adjusting the base-pair
scale.  The resulting conditional bounds are
\[
 \RFE T(N)\le L_N[1/3,91.4],\qquad S(N)\le L_N[1/3,19].
\]
The conditional qualification is part of the theorem: the weighted
progression estimates, modular polynomial coupling, global descent
readiness, and algebraic square-certification assertions are specified
as a proof contract in the enlarged parameter range.  The source's
fixed-parameter statements alone do not establish that contract.  We
prove the sparse kernel and hidden-root statements independently,
derive the revised recurrence, and supply explicit checks of the
changed parameter inequalities.  An initially proposed schedule fails
a root progression-modulus condition; the repaired schedule uses
$\tau=19/16$ and $g=43/8$ and is recorded with that failure visible.
An exact rational verifier certifies the finite coefficient inequalities,
but does not certify the analytic contract.  The first $L$-parameter
remains $1/3$; no unconditional new factoring record, practical running
time claim, or bound with first parameter below $1/3$ is asserted.
\end{abstract}

\tableofcontents

\section{Scope, notation, and comparison}
\label{rf:scope}

The goal of this manuscript is to make a proposed optimization
reconstructible and to state exactly which parts are proved here.
The underlying factoring construction is the version of
\cite{CogRand2026} at Google Research commit
\texttt{4a7bc8ab0bc770e71a259b7966c500496b99010a}; the PDF has Git blob
\texttt{19f36b2f189ed41ec8e112b92aa9b88a96c6bd1e} and 54 pages.
Source page numbers below refer to that PDF, rather than to an evolving
repository version.  The separate provenance record and source-fetch
script identify the same bytes.

For $N>e^e$, put
\begin{equation}
 L=\log N,\quad \ell=\log L,\quad
 U=(L/\ell)^{1/3},\quad
 \Lambda=L^{1/3}\ell^{2/3},\quad
 \Xi=L^{2/3}\ell^{1/3}.
 \label{rf:scales}
\end{equation}
Thus $\Lambda=U\ell$, $\Xi=U\Lambda$, and $L=U\Xi$.
We use the source's distinction between
\[
 L_N(\alpha,c)=\exp\{c(\log N)^\alpha
                         (\log\log N)^{1-\alpha}\}
\]
and the assertion $T(N)\le L_N[\alpha,c]$, which means that the
corresponding normalized logarithmic limsup is at most $c$.
Writing $\exp((c+o(1))\Lambda)$ always concerns a uniform asymptotic
upper bound over the inputs and internal states under discussion.
Constants in a proof contract must be effective when effectivity is
used in the algorithmic conclusion.

The machine model is a uniform classical probabilistic multitape
Turing machine.  Time counts bit operations.  Workspace includes all
active recursive tapes and is bounded for every coin tape.  Sparse
matrix operations below are implemented by sequential sorting and
merging; constant-time random access to an exponentially long array
is not assumed.

\begin{center}
\begin{tabular}{p{.38\textwidth}ccp{.24\textwidth}}
\toprule
Construction or assertion & Time & Space & Status in this manuscript\\
\midrule
Pinned source, Theorem 1.1, p.~2
 & $L_N[1/3,105]$ & $L_N[1/3,40]$ & Source theorem claim\\
Pinned source, induction, pp.~47--49
 & $L_N[1/3,104]$ & $L_N[1/3,38]$ & Stronger printed accounting\\
Sparse implementation with the old root-prime scale
 & $L_N[1/3,102]$ & $L_N[1/3,19]$ & Conditional precursor\\
Repaired enlarged-cap schedule studied here
 & $L_N[1/3,91.4]$ & $L_N[1/3,19]$ & Conditional on the stated contract\\
\bottomrule
\end{tabular}
\end{center}

This table compares the present construction with the pinned source;
it is not a certification of a world record.  Lenstra and Pomerance
proved a classical expected-time bound $L_N[1/2,1]$
\cite{LenstraPomerance1992}.  The familiar
$L_N[1/3,(64/9)^{1/3}]$ analysis of the general number field sieve uses
heuristic assumptions \cite{GNFS1993}.  The much larger second
constants in the present conditional construction therefore should
not be interpreted as a practical improvement on number field sieve
implementations.

\subsection{What is changed, and what is retained}

There are three changes with different purposes.
First, a sparse, verifiable kernel routine replaces dense Gaussian
elimination.  This changes the leading workspace accounting.
Second, rejected trial roots are erased, while accepted roots remain
sealed until a square relation has been selected.  This removes an
unnecessary array of all trial roots.
Third, the accepted top-level numerator and denominator are products
of fewer, larger primes.  Their occurrence probability improves,
reducing the number of recursive factorizations of integers near
$\sqrt N$.

The last change is not free.  It enlarges some progression moduli and
some descendant norms.  Increasing only the prime cap fails; the
base-pair scale must also increase.  Section~\ref{rf:rootrepair}
derives this failure and its repair.  The polynomial degree,
coefficient-height constant, factor-base constant, two-field descent,
exact square-root tests, and final congruence-of-squares verification
are retained.

\section{The conditional theorem and its proof contract}
\label{rf:contract}

We state the main conclusion before describing the construction in
detail.  The contract below is operational: it identifies the exact
uniform success and certification statements used in the complexity
proof.  Section~\ref{rf:analyticcontract} states the corresponding
analytic and coupling obligations more explicitly.

\begin{assumption}[Enlarged-parameter source contract]
\label{rf:assumption}
There is an effective threshold $N_0$ such that the following
assertions hold for every odd composite, non-perfect-power integer
$N\ge N_0$ having no prime divisor at most $B$, with the parameters
and finite procedures in Sections~\ref{rf:parameters}--\ref{rf:algorithm}.
\begin{enumerate}[label=\textup{(C\arabic*)},leftmargin=*]
\item \textbf{Setup and uniform global readiness.}
The capped base-pair and polynomial setup succeeds with probability
at least $3/4$ and produces a tuple
$(g_0,g_1,f_1,f_2)$ satisfying all of its stated exact algebraic checks.
For each such good tuple there is a set
\[
 \mathcal Q_{\rm good}\subseteq
 \{q\text{ prime}:Q_0/2<q\le Q_0,
                   \ q\nmid 3Ng_0g_1c_{1,d}c_{2,d}\}
\]
with at least $3/4$ as many elements as
$\mathcal P_0=\{q\text{ prime}:Q_0/2<q\le Q_0,
q\nmid3Ng_0g_1\}$.
For every fixed $q\in\mathcal Q_{\rm good}$, with the tuple held fixed,
the capped procedure $\operatorname{DescentTree}(q)$ returns a valid
tree or a certified proper factor of $N$ with probability at least
$1-L^{-2}$.  This probability is over fresh tree coins.  The sets of
good primes are analytical objects and need not be computed.

\item \textbf{Algebraic relation bookkeeping.}
Every accepted tree has the exact signed relation described in
Section~\ref{rf:relations}.  After all degree-one parity rows and
sampled character rows vanish, its class lies in the joint residual
square-class obstruction space of the two fields.  The local length
formulas and the enlarged obstruction groups of source Lemma~5.1
apply to the resulting lists.  They apply also at ramified primes
in the polynomial order, in the precise length sense used there.

\item \textbf{Character certification.}
For each tuple from \textup{(C1)}, the capped character sampler of
source Lemma~5.1 succeeds and leaves a joint residual obstruction
space of dimension at most two with probability at least $98/100$.
Its time and deterministic space are at most
$\exp((8+o(1))\Lambda)$ and $\exp((4+o(1))\Lambda)$,
respectively.  Its coins are independent of the future sealed roots.

\item \textbf{Exact algebraic square extraction.}
For the forms and positive factor lists used below, source
Lemma~5.2 is valid as a deterministic decision and extraction
algorithm.  With
\[
 B_{\rm bit}=\left\lceil
 4d^3\log_2(2dM)+M_{\rm fac}\log_2(2N)+d\log_2(2M)
 \right\rceil,
\]
it decides the required field-square predicate and, when true,
returns the integral polynomial square root in
$O(d^3B_{\rm bit}^2)$ bit operations and
$O(d^2B_{\rm bit})$ space.  Final roots and modular squares are
checked exactly.

\item \textbf{Uniformity of the source procedures.}
The capped samplers, descent filters, weight evaluation, and
bookkeeping have the finite bit-operation bounds specified below;
there is no uncharged factoring oracle, unbounded precision oracle,
or unbounded-memory random loop.  All errors in \textup{(C1)}--\textup{(C4)}
and the underlying analytic statements are uniform in their stated
parameters.  The range changes in
Section~\ref{rf:analyticcontract} are part of this assumption.
\end{enumerate}
\end{assumption}

The principal unproved part of this contract in the present manuscript
is the independent certification of the source's analytic, conditional
moment, and global-readiness arguments in the enlarged parameter
range.  The contract is not an assumption that a factor is known.
Every factor returned by the machine is verified without it.  Its
role is to justify the promised expected running time.

\begin{theorem}[Conditional enlarged-cap factoring bound]
\label{rf:main}
Assume Assumption~\ref{rf:assumption}.  The modified machine described
below completely factors every integer $N\ge2$, with verified prime
outputs, and has
\[
 \RFE T(N)\le L_N[1/3,91.4],\qquad
 S(N)\le L_N[1/3,19].
\]
The time bound is worst-case expected time over input integers of
the given size.  The space bound holds on every coin tape.  The
machine retains deterministic background trial division and hence
terminates on every coin tape, including inputs below $N_0$.
\end{theorem}

The conditional nature of the theorem is essential.  The source
lemmas are stated using $\tau=1/3$ and $g=3$.  They cannot simply be
substituted as black boxes after those constants change.  Our
parameter audit derives the replacements needed in their proofs,
including one condition that the first proposed schedule violated;
it does not purport to certify every line of the source's long
analytic argument.

\section{Parameters, exact data, and the source interface}
\label{rf:parameters}

\subsection{The modified schedule}

Use
\begin{equation}
 \delta=\frac12,\quad \tau=\frac{19}{16},\quad
 g=\frac{43}{8},\quad \mu=8,\quad \beta=18,\quad
 \eta_0=2,\quad \rho=\frac34.
 \label{rf:constants}
\end{equation}
Dyadic approximations are taken as in source Definition~2.7:
the coarse logarithmic scales have fixed absolute error at most
$2^{-20}$, and fine logarithms and exponentials use precision
$P_{\log}=\lceil60\widehat\ell^{,2}+40\rceil$.
Rounding the fine approximations gives errors at most
$2^{-P_{\log}}$ in logarithms and relative errors of that order in
exponentials.  All integer formulas below mean these deterministic
dyadic implementations, rather than access to exact real numbers.

Set
\begin{align}
 d&=2\left\lfloor\frac{\delta}{2}\widehat U\right\rfloor+1,
 &r_0&=\left\lceil\frac{\widehat U}{2\tau}\right\rceil,
 &Q_0&=\max\{q\in\RFZ_{\ge1}:4q^{2r_0}\le N\},
 \label{rf:dqr}\\
 B&=\lceil\exp(\beta\widehat\Lambda)\rceil,
 &G&=\lceil\exp(g\widehat\Xi)\rceil,
 &M&=3\lceil\exp(\mu\widehat\Xi)\rceil.
 \label{rf:bgm}
\end{align}
In particular,
\begin{equation}
 d=(1/2+o(1))U,\quad
 r_0=(8/19+o(1))U,\quad
 \log Q_0=(19/16+o(1))\Xi,
 \quad \frac N8\le Q_0^{2r_0}\le\frac N4
 \label{rf:rootasymptotics}
\end{equation}
for sufficiently large $N$.  The lower bound follows from
$Q_0>(N/4)^{1/(2r_0)}-1$ and $r_0/Q_0\longrightarrow0$.

Retain
\begin{align*}
 K_\xi&=\lceil\exp(30\widehat\ell^{,2})\rceil,
 &W&=\lfloor M/(6K_\xi)\rfloor,
 &W_{\rm cut}&=\lceil W\exp(-\widehat\Xi)\rceil,\\
 \Lambda_0&=\lceil32\widehat\Lambda\rceil,
 &S_{\rm skew}&=\lceil\exp(\widehat\Lambda/2)\rceil,
 &M_{\rm bins}&=\lceil\exp(19\widehat\Lambda/100)\rceil,\\
 M_{\rm wed}&=\lceil\exp(45\widehat\ell^{,2})\rceil,
 &H_0&=\lceil\exp(21\widehat\Lambda/5)\rceil,
 &K_0&=\lceil\exp(160\widehat\ell^{,2})\rceil,\\
 E_{\rm gcd}&=\lceil\exp(\widehat L^{,1/4})\rceil.
\end{align*}
For integer $q>B$, let
\[
 E_q=8\left\lceil\frac18\sqrt q\exp(2\widehat\Lambda)\right\rceil,
 \qquad
 \mathcal B_q=
 \bigl([3E_q/4,E_q]\times[E_q/8,E_q/4]\bigr)\cap\RFZ^2.
\]
For sufficiently large $N$,
$\log E_q=\tfrac12\log q+2\Lambda+o(1)$.
The factor eight makes the box endpoints integral.

Define the geometric ladder by
\begin{gather*}
 J=\min\{j\ge1:\rho^j\widehat{\log Q_0}
                       \le\widehat{\log B}+1\},\\
 Y_0=Q_0,\quad
 Y_j=\lfloor\exp(\rho^j\widehat{\log Q_0})\rfloor\ (0<j<J),
 \quad Y_J=B.
\end{gather*}
The bucket $I_j=(Y_{j+1},Y_j]$ has child cap
$y(j)=Y_{\min(j+2,J)}$.  Children move at least two buckets downward.
For fixed constants as in \eqref{rf:constants},
\begin{equation}
 J\le1.2\ell,\quad
 r_{\max}\le\lceil J/2\rceil\le0.65\ell,\quad
 \frac{\log q}{\log y(j)}\le\frac{16}{9}+10^{-3},\quad
 E_q\le q\exp(-6.9\Lambda)
 \label{rf:ladderbounds}
\end{equation}
eventually.  Rounding changes neither these strict inequalities nor
their effective eventual validity.

Two analytical quantities used only in proofs are
\begin{equation}
 M_{\rm pair}=\lceil\exp(0.19\Lambda)\rceil,
 \qquad M_G=\lceil\exp(30\ell^2)\rceil.
 \label{rf:analyticaux}
\end{equation}
The source used $M_{\rm pair}=\exp(0.15\Lambda)$.
Changing $M_{\rm pair}$ does not add an algorithmic search loop.
Similarly, exceptional conductor sets introduced later are not
computed by the algorithm.

\subsection{Tuned smoothness parameters and exact weights}

Let $u_{\rm cap}=\lceil4\widehat U\rceil$,
$L_{\rm grid}=\lceil6u_{\rm cap}\widehat{\log Q_0}\rceil$, and
$\kappa_0=(\log2)/2$.  For $2\le k\le J$ and
$0\le v\le L_{\rm grid}$, predefine on demand
\[
 y_{k,v}=\left\lfloor
      Y_k\exp\{-v/(10u_{\rm cap})\}\right\rfloor.
\]
For a positive rational $X\in[2,N]$, set
\begin{align}
 u(X,Y_k)&=
 \left\lceil\frac{\widehat{\log X}}
                   {4(\widehat{\log Y_k}-1/2)}\right\rceil,
 \label{rf:tune-u}\\
 v(X,Y_k)&=
 \left\lfloor10u_{\rm cap}
 \left(\widehat{\log Y_k}-\widehat\kappa_0-
       \frac{\widehat{\log X}}{4u(X,Y_k)}\right)\right\rfloor,
 \label{rf:tune-y}
\end{align}
and $y(X,Y_k)=y_{k,v(X,Y_k)}$.
Out-of-range values use the source sentinel $(y,u)=(0,10)$ and are
rejected.  All factorization calls below precede evaluation of the
weight; there is no inexpensive smoothness oracle hidden here.

For an integer $y\ge2^{100}$, let
$\mathcal P(y)=\{p\text{ prime}:y/2<p\le y\}$ and
$\mathcal I(y)=(y/2,y]\cap\RFZ$.  Define
\begin{equation}
 \widetilde w_{y,u}(n)=\frac1{(4u)!}
 \#\left\{(m,p_1,\ldots,p_{4u-1}):
 \begin{array}{l}
 m\in\mathcal I(y),\ p_i\in\mathcal P(y),\\
 n=mp_1\cdots p_{4u-1}
 \end{array}\right\},
 \label{rf:weight}
\end{equation}
and set it to zero for $n\le0$.  It lies in $[0,1]$ and its support
consists of $y$-smooth integers in $((y/2)^{4u},y^{4u}]$.
For clarity, its exact evaluation requires only the complete prime
factorization of $n$.  Write
\[
 n=M_{\rm sm}\prod_{p\in\mathcal P(y)}p^{e_p},
 \quad P^+(M_{\rm sm})\le y/2,\quad E=\sum e_p.
\]
If $n$ has a prime factor greater than $y$, the weight is zero.
Otherwise it is
\begin{equation}
 \widetilde w_{y,u}(n)=
 \begin{cases}
 (\prod e_p!)^{-1},&M_{\rm sm}=1,\ E=4u,\\
 (4u\prod e_p!)^{-1},&M_{\rm sm}\in\mathcal I(y),\ E=4u-1,\\
 0,&\text{otherwise}.
 \end{cases}
 \label{rf:weightformula}
\end{equation}
Indeed the free integer $m$ either is the one remaining large prime
or is exactly $M_{\rm sm}$.  Thus a rational Bernoulli acceptance
probability can be evaluated and sampled exactly with capped uniform
integer draws.

\subsection{The polynomial ensemble}

Choose a base pair uniformly from
\[
 \mathcal G=\{(g_0,g_1)\in([G,2G]\cap\RFZ)^2:
       \gcd(g_0,g_1)=\gcd(g_0g_1,3N)=1\}.
\]
Choose a monic irreducible $h_0\in\RFF_3[X]$ of the odd degree $d$.
Let $m_N\equiv g_0g_1^{-1}\pmod N$ and pin the coefficients modulo
three as in source Definition~4.1.  Explicitly, if
$h_0=X^d+\sum_{i<d}\bar h_iX^i$, put
\begin{gather*}
 \bar c=(\bar h_0,\ldots,\bar h_{d-1},M+3W+1),\qquad
 c(\widetilde c)=\bar c+3\widetilde c,\\
 f_{\widetilde c}(X,Y)=\sum_{i=0}^d c_i(\widetilde c)X^iY^{d-i}.
\end{gather*}
The affine lattice condition is
\begin{gather*}
 \widetilde{\mathcal L}_N=
 \left\{\widetilde c\in\RFZ^{d+1}:
       \sum_{i=0}^d\widetilde c_i m_N^i\equiv r_N\pmod N\right\},\\
 3g_1^dr_N\equiv-f_0(g_0,g_1)\pmod N.
\end{gather*}
On $\widetilde{\mathcal B}=[-W,W]^{d+1}\cap\RFZ^{d+1}$ use the
weight
\[
 \exp\{-\Lambda_0\|\widetilde c\|_1/W\}.
\]
Its normalized law before the lattice condition is $P_0$; the law
conditioned on $\widetilde{\mathcal L}_N$ is $P_{\rm Lap}$.
The source's finite layer-cake/LLL sampler, with the stated dyadic
precision and rejection caps, is used to sample two independent
forms.  It expresses the product Laplace law as a mixture of centered
integer boxes, samples the box radii, and samples lattice points in a
box by tiling with an LLL fundamental parallelepiped.  Its finite
sampling law $P_{\rm alg}$ is within $\RFLT{-1/2}$ of $P_{\rm Lap}$
in total variation under the contract.  LLL itself is a deterministic
polynomial-time lattice algorithm \cite{LLL1982}.

For every sampled coefficient vector, the following are exact
algebraic properties or direct size estimates:
\[
 f(g_0,g_1)\equiv0\pmod N,\quad
 f(X,1)\bmod3=h_0(X),\quad
 c_d\in[M+1,M(1+2K_\xi^{-1})],
\]
and for $(a,b)\in\mathcal B_q$,
\begin{equation}
 Ma^d(1-K_\xi^{-1})\le f(a,b)
                   \le Ma^d(1+2K_\xi^{-1}).
 \label{rf:normsize}
\end{equation}
Reduction modulo three proves irreducibility over $\RFQ$ and excludes
three from values at primitive pairs.

\subsection{Finite guards and all-coin correctness}

Retain the source's exact guards
\begin{gather*}
 d\ge9,\quad r_0\ge10,\quad J\ge4,\quad
 2^{100}\le B=Y_J<2B-1<Y_{J-1}<\cdots<Y_0,\\
 d>J+4,\qquad G^{d-J-4}>N,\qquad W\ge8W_{\rm cut},\\
 W_{\rm cut}\ge16(d+1)
       \left\lceil(d+1)2^{(d+3)/2}G\right\rceil.
\end{gather*}
If they fail, use deterministic trial division.  Every recursive
call inside a randomized phase is additionally guarded by the exact
range $2\le m\le\lfloor\sqrt N/2\rfloor$.  An out-of-range candidate
is rejected.  A proper gcd is verified before it is returned.
Primality tests and perfect-power extraction are deterministic;
AKS provides a suitable certified primality test \cite{AKS2004}.
These conditions guarantee correctness independently of any density
claim.  The asymptotic contract guarantees that the fallback and
range rejections do not spoil the claimed expected bound above an
effective threshold.

\section{The modified algorithm}
\label{rf:algorithm}

\subsection{Complete factorization and a single phase}

The full procedure maintains a worklist of unresolved factors.
It strips powers of two, extracts perfect powers, certifies primes,
and trial-divides unresolved non-perfect-power composites by all
integers up to $B$.  It calls the following splitting procedure only
on the remaining odd composites.  Each returned proper divisor is
checked before the worklist is updated.  On completion, sorted prime
powers are certified and their product is checked against $N$.

\begin{algorithm}[A capped randomized splitting phase]
\label{rf:phasealgorithm}
For a guarded input $N$, run the following phase.  Repeat fresh
phases after a failure, overwriting their scratch space.  Interleave
every transition with one transition of deterministic trial division
on $N$; either thread may finish the split.
\begin{enumerate}[leftmargin=*]
\item Sample $(g_0,g_1)$ and two forms from the ensemble above.
Check the leading-coefficient and derivative gcds required for
modular evaluation.  Return any proper gcd; restart on an invalid
tuple or sampling timeout.
\item Enumerate the rational primes and degree-one polynomial-order
prime ideals over $p\le B$.  Sample the character pairs of source
Lemma~5.1.  Let
\[
 D_{\rm row}=3+\pi(B)+|\mathcal F_1(B)|+|\mathcal F_2(B)|+2L_{\rm ch},
 \quad R_{
 \rm tot}=D_{\rm row}+3,
 \quad L_{\rm ch}=4\lceil(\widehat L+1)^2\rceil+40.
\]
Only sparse relation records are allocated.
\item Put $\epsilon_N=\ell^{-1/2}$, using a positive dyadic
approximation of error $o(\epsilon_N)$, and set
\begin{equation}
 M_{\rm cap}=\left\lceil8R_{\rm tot}
       \exp\left((16/19+\epsilon_N)\widehat\Lambda\right)\right\rceil.
 \label{rf:trialcap}
\end{equation}
Attempt at most $M_{\rm cap}$ independent blind trials, stopping
when $R_{\rm tot}$ relations have been accepted.
\begin{enumerate}[label=\textup{(\alph*)},leftmargin=*]
\item Draw $z$ uniformly from $\{1,\ldots,N-1\}$.
A nontrivial $\gcd(z,N)$ is returned immediately.  Otherwise seal
$z$ and expose only $w=z^2\bmod N$.
\item By Lagrange--Gauss reduction, compute a shortest nonzero
vector $(x,y)$ of
$\{(x,y)\in\RFZ^2:x\equiv wy\pmod N\}$, with the deterministic
sign convention $x>0$.  Check units and
$2\le x,y\le\lfloor\sqrt N/2\rfloor$.  Call the full factorer on
both integers.  Require that they are products of disjoint
$r_0$-element sets of distinct primes in $\mathcal P_0$, none
dividing $c_{1,d}c_{2,d}$.
\item Descend each of these $2r_0$ primes by
$\operatorname{DescentTree}$.  A verified split terminates the
phase successfully; any failure rejects the trial.  If all trees
succeed, assemble the signed relation and its sparse parity column.
Retain its sealed root.  Erase the root of every rejected trial.
\end{enumerate}
\item If too few relations were collected, restart.  Otherwise run
at most three independent attempts of the sparse kernel routine in
Section~\ref{rf:sparse}, obtaining three verified independent binary
null vectors or restarting on failure.
\item Test the seven nonzero combinations of those vectors by the
exact algebraic square-root routine.  Select the first combination
that passes both field-square tests, construct the modular root
$Y$, and verify that it is a unit and
$Y^2=\prod_{m\in T}w_m\pmod N$.
\item Only now unseal the accepted roots indexed by $T$, put
$Z=\prod_{m\in T}z_m\pmod N$, and test $\gcd(Z-Y,N)$.
Return it if proper; otherwise restart.
\end{enumerate}
\end{algorithm}

All elementary uniform draws have the source cap
$\lceil10(\widehat L+1)^2\rceil$ attempts.  The other source
sampler caps are retained: polynomial/box draws have polynomial
caps, character outer proposals have cap
$\lceil\exp(2\widehat\Lambda)\rceil$, and finite-field polynomial
splitting uses the prescribed polynomial cap.  The sparse kernel
uses exact uniform bit strings and finite loops.  The result is a
single specified machine with an all-coin space bound, rather than
a procedure that assumes every randomized subroutine eventually
finishes within its expected budget.

\subsection{Root and algebraic descent steps}
\label{rf:descent}

For reproducibility, the descent interface is described here rather
than hidden inside the time recurrence.  A prime $q$ and a residue
$r$ define
$\Lambda_{q,r}=\{(a,b):a\equiv br\pmod q\}$.
It is non-skewed when its shortest nonzero vector has Euclidean
length at least $\sqrt q/S_{\rm skew}$.  The source's exact lattice
box sampler draws uniformly from $\Lambda_{q,r}\cap\mathcal B_q$,
with success per proposal at least $1/825$ when non-skewed.

At a root prime $q_0$, put $r=g_0g_1^{-1}\pmod{q_0}$ and use
\[
 T_{\rm try,0}=\left\lceil
 8(\widehat L+1)^3\left\lceil\frac{E_{q_0}^2}{31q_0}\right\rceil
       \exp(-2.75\widehat\Lambda)\right\rceil.
\]
For each proposal, require $(a,b)$ primitive, check the gcd of
$ag_1-bg_0$ with $N$, and form
\[
 n_{\rm rat}=(ag_1-bg_0)/q_0,\qquad
 n_{
 \rm alg}=f_1(a,b),\qquad n'_{
 \rm alg}=f_2(a,b).
\]
Check all three size guards.  Factor $n_{\rm rat}$ and $n_{\rm alg}$.
Require the rational weight, tuned to $(n_{\rm rat},B)$, to be
positive and its primes in $\mathcal P(y)$ to have valuation at most
one.  Accept with exact probability equal to the algebraic weight
tuned to $(Ma^d,Y_2)$.  After acceptance, factor $n'_{\rm alg}$ for
the recorded history, and recurse on the primes $p>B$ occurring
with odd valuation in $n_{\rm alg}$, omitting primes dividing
$c_{1,d}$ and enforcing the source unit and residue guards.

An algebraic node has data $(\iota,q,r;\mathcal H)$, where
$\iota\in\{1,2\}$, $\iota'=3-\iota$, and the history records the
ancestor pairs and both of their form values.  If $q\in I_j$, put
\[
 T_{\rm try}(q)=\left\lceil
 8(\widehat L+1)^3\left\lceil\frac{E_q^2}{31q}\right\rceil
             \exp(-2\widehat\Lambda)\right\rceil.
\]
Check non-skewness, all required cross-value nondivisibilities,
and that $q$ does not divide any determinant between two ancestors.
For each sampled primitive pair $\boldsymbol a=(a,b)$, let
\[
 \Delta_t=|b_t^*a-a_t^*b|,\qquad
 \Delta'_t=\Delta_t/q^{v_q(\Delta_t)}.
\]
Check the source bounds $1\le\Delta_t\le\sqrt N/2$ and
$M_{\rm wed}\Delta_t\ge a_t^*E_q$, factor every $\Delta'_t$, and
perform the following exact pruning tests:
\begin{enumerate}[label=\textup{(\roman*)},leftmargin=*]
\item $\gcd(\Delta'_t,y_{t,1}^*y_{t,2}^*)\le E_{\rm gcd}$;
\item the product of primes dividing both $\Delta'_t$ and an
ancestor determinant is at most $E_{\rm gcd}$;
\item
$\prod_{p\mid\Delta'_t,\ p^2\le H_0,\ v_p(\Delta'_t)\ge2}
 p^{\min(v_p(\Delta'_t),e_{H_0}(p))}\le K_0$,
where $e_{H_0}(p)=\max\{e:p^e\le H_0\}$.
\end{enumerate}
Form $C_\iota=f_\iota(a,b)/q$ and
$C_{\iota'}=f_{\iota'}(a,b)$, check the exact ranges, and factor
both.  Accept with probability
\[
 \widetilde w_{y_1,u_1}(C_\iota)
 \widetilde w_{y_2,u_2}(C_{\iota'}),
\]
where the parameters are tuned to $(Ma^d/q,y(j))$ and
$(Ma^d,y(j))$.  Record the relation and recurse on each required
odd-valuation prime in $(B,y(j)]$, with the enlarged history.
A failed guard, trial exhaustion, or child failure gives the specified
failure outcome.  This is source Definition~4.5 with the new global
parameters, not a replacement of recursive factorization by a
heuristic smoothness test.

\begin{proposition}[Deterministic descendant bounds]
\label{rf:treebounds}
For the modified parameters, every guarded descent has at most
$\exp(O(\ell^2))$ nodes.  Uniformly over its prospective evaluated
integers,
\[
 |f_\iota(a,b)|\le N^{19/64+o(1)},\qquad
 \log n_{\rm rat}\le(153/32+o(1))\Xi,\qquad
 \log\Delta_t=O(\Xi).
\]
For sufficiently large $N$, all are below $\sqrt N/2$.
\end{proposition}
\begin{proof}
The ladder gives depth at most $0.65\ell$.  By
\eqref{rf:normsize},
\[
 \log|f_\iota(a,b)|\le\log(2M)+d\log E_{Q_0}
                    =(\delta\tau/2+o(1))L
                    =(19/64+o(1))L.
\]
The product of two algebraic cofactors has exponent at most
$19/32+o(1)<1$, so the number of distinct potential children is
at most $\log_2N$; exact range guards also preserve this crude
bound on all finite inputs.  Summing the resulting geometric tree
bound gives $\exp(O(\ell^2))$ nodes.  The rational estimate follows
from $n_{\rm rat}\le4E_{Q_0}G/Q_0$, and determinants are bounded
by products of two special-box scales.  Since $\Xi=o(L)$ and
$19/64<1/2$, the asserted size guards eventually hold.
\end{proof}

\section{Analytic obligations and parameter transfer}
\label{rf:analyticcontract}

This section states which source ranges must be widened and proves
the numerical comparisons used in that proposed transfer.  The
logical status is deliberate: the calculations are proofs of
parameter inequalities.  They are not a substitute for an
independent proof of the analytic and conditional-probability
statements listed next.

\subsection{The weighted progression obligation}

For $y\in[B^{0.99},Q_0]$, call $(y,u)$ admissible when
\[
 \max(10,U/30)\le u\le u_{\rm cap},\qquad
 4u\le0.96U,\qquad (8u\log y)^u\le\exp(0.25\Lambda).
\]
Put $Q_{\rm AP}(y,u)=\lfloor y^{u/2-1/16}\rfloor$ and
$\theta_{\rm ls}=\exp(-0.6\Lambda)$,
$T_{\rm Mel}=\exp(\Lambda)$.
The enlarged version of source Lemma~3.5 required here has the
same formulas and error bounds as its original version, but must
hold for this enlarged $Q_0$ and the full tuned grid.
It provides exceptional conductor sets
$\mathcal F_{\rm bad}(y,u)\subset(y^{3/5},Q_{\rm AP}]$ of size at
most $\exp(1.46\Lambda)$ and a density
\begin{equation}
 \nu(X;y,u)=\frac1{(4u)!}
 \sum_{p_1,\ldots,p_{4u-1}\in\mathcal P(y)}
 \frac{\boldsymbol1_{X/y\le P\le2X/y}}P,
 \qquad P=p_1\cdots p_{4u-1}.
 \label{rf:density}
\end{equation}
For $|4u(\log y-\kappa_0)-\log X|\le2$, it requires in particular
\begin{equation}
 \nu(X;y,u)\ge\tfrac12(5u\log y)^{-4u}
                  \ge\exp(-0.97\Lambda),
 \label{rf:densitylower}
\end{equation}
the source's short-scale continuity estimate, and its polynomial
comparability for nearby tuned scales.

More precisely, let $z_0,R,m\ge1$ be integers,
$g_m=\gcd(z_0,m)\le y^{3/10}$,
$\omega(m)\le2\log N$, and
$|4u(\log y-\kappa_0)-\log z_0|\le2$.
Write $H=2mR$ and $X_0=z_0-mR$.  When
\begin{equation}
 m\le Q_{\rm AP},\quad
 f\nmid m\quad(f\in\mathcal F_{\rm bad}),\quad
 X_0e^{-0.2\Lambda}\le H\le X_0e^{-29\ell^2},
 \label{rf:apconditions}
\end{equation}
the required conclusion is
\begin{equation}
 \left|\frac1{2R+1}\sum_{t=-R}^R
       \widetilde w_{y,u}(z_0+mt)-\nu(z_0;y,u)\right|
 \le e^{-27\ell^2}\nu(z_0;y,u).
 \label{rf:apclaim}
\end{equation}
The unconditional upper-bound portion of the same obligation is
\begin{equation}
 \frac1{2R+1}\sum_{t=-R}^R\widetilde w_{y,u}(z_0+mt)
 \le20\sqrt u\,(u'+1)\{3(u'+1)\log y\}^{u'}\nu(z_0;y,u)
 \label{rf:roughclaim}
\end{equation}
for $1\le u'\le2u$, $H\le z_0/2$, and
$(y/2)^{u'+1}\ge6m(z_0/H)$, without the good-conductor
condition.  The associated bad-child estimate from source
Lemma~3.6(3) has the factor
$\{3(u'+1)\log y\}^{u'}\exp(-0.93\Lambda)$ and the slightly
stronger prefix condition $(y/2)^{u'}\ge6m(z_0/H)$.

The exceptional sets are united over the grid to form
$\mathcal S_{\rm bad}$, and primes $q>B$ dividing $N$ or one of
its members form $\mathcal Q_{\rm exc}$.  The needed cardinalities
are $\exp(1.48\Lambda)$ and $\exp(1.49\Lambda)$.
They are proof devices, not extra input to the finite descent.

\subsection{The coupling and conditional-moment obligations}

The enlarged version of source Lemma~4.2 must allow
$g=43/8$ and evaluation vectors
$\boldsymbol v_1,\ldots,\boldsymbol v_R$ with
\[
 R\le\lceil J/2\rceil+3,\quad
 v_{s,1}\ge3v_{s,2}\ge3,\quad
 \prod_s(2v_{s,1})\le\exp((34/25)\Xi),\quad
 \gcd(v_{s,1}g_1-v_{s,2}g_0,N)=1.
\]
The vectors are primitive and pairwise non-collinear.  The slice
may also impose $f(r_q,1)=0\pmod q$ for $B<q\le Q_0$,
$q\nmid3N$.  With
$\widetilde{\mathcal B}_{\rm core}
=\{\|\widetilde c\|_1\le W/4\}$ and
$\widetilde{\mathcal B}_{\rm mid}
=\{\|\widetilde c\|_1\le W/2\}$,
the required conclusions are the original fibre tail and transfer
bounds: a slice meeting the core has tail outside the middle box
at most $e^{-5\Lambda}$, and any $[0,1]$-valued function of the
listed evaluations changes expectation under the mod-$N$
conditioning by at most
\[
 2e^{-5\Lambda}\RFE_{P_0}(\Psi)+5e^{-5\Lambda}.
\]
The source's multiplicative Laplace sub-averaging bound is retained
for shifts satisfying its norm and tail hypotheses.

Source Lemma~4.4 must then hold for all valid histories of the new
ladder, rather than only for the original prime cap.  A history
contains primitive ancestor pairs at levels separated by at least
two buckets, their two form values, a current non-skewed prime
$q\notin\mathcal Q_{\rm exc}$, and the stated nondivisibilities by
ancestor determinants and opposite-field values.  Its conditional
polynomial slices must meet the core.  For the source's clean set
of candidate pairs, the required local conclusions are
\begin{gather*}
 \mu_{\rm node}\ge e^{2.04\Lambda},\qquad
 \RFPr\{Z_{\rm clean}<\mu_{\rm node}/2\}\le e^{-25\ell^2},\\
 \RFE X_{\boldsymbol a}\le
 (1+e^{-26\ell^2})\nu_1(\boldsymbol a)\nu_2(\boldsymbol a),\qquad
 \RFE\sum X_{\boldsymbol a}\boldsymbol1_{\rm bad\ child}
       \le e^{-0.4\Lambda}\mu_{\rm node}.
\end{gather*}
At a rational root the corresponding weight is the single
algebraic weight, the mean is at least $e^{2.80\Lambda}$, and the
same concentration and bad-child conclusions are required.
Source Lemma~4.6's backwards conditional-expectation argument must
apply to these enlarged slices, giving \textup{(C1)}.  It is this
uniform conditional statement, not merely an average smoothness
probability, that is needed for adaptive descent.

\begin{remark}[Limits of the transfer audit]
The source's Fourier, large-sieve, local-limit, lattice-conditioning,
and backwards-telescoping arguments have not all been independently
reproved here.  The displays above state the actual enlarged
assertions required.  The following sections verify the changed
size inequalities and locate the original constants that must be
replaced.  They do not turn the assertions into unconditional
theorems by referring to a fixed-parameter source statement.
\end{remark}

\subsection{Grid size and universal admissibility}

The number of grid index pairs is now bounded by $40L\ell$ for
sufficiently large $N$, in place of the source's $12L\ell$.
Including $u$ gives at most $200L^2$ admissible triples.
Polynomial unions therefore preserve the exceptional-set
coefficients $1.48$ and $1.49$.
For actual algebraic-node scales, the source tuning calculation gives
\[
 \frac{4u}{U}\le
 \frac{\mu+\delta\eta_0}{\beta}+
       \frac{\delta}{2\rho^2}+o(1)
 =\frac{17}{18}+o(1)<0.96.
\]
At every algebraic node, $\log y\le\log q$ and
$\log(Ma^d/q)\ge(d/2-1)\log q+O(\Xi)$ give the lower bound
\begin{equation}
 2u\ge(1/8+o(1))U.
 \label{rf:algu-lower}
\end{equation}
This lower bound is stronger than the universal admissibility bound
and is used in the close-pair estimate below.

In the general admissible domain, $u\le0.24U$ and
$\log y=O(\Xi)$ imply
$\log(8u\log y)\le\ell+O(1)$; hence the printed
$\exp(0.25\Lambda)$ bound is retained.
The source's dispensable inequality $5u\log y<L$ need not remain
literal throughout the enlarged abstract domain.  Replacing it by
$5u\log y\le C L$ changes the density logarithm only by $O(U)$,
which is $o(\Lambda)$, and preserves \eqref{rf:densitylower}.
Also $\log(y^u)\le(0.24\tau+o(1))L=(0.285+o(1))L<L$;
the large-sieve support and conductor prefactors bounded using
$\log N$ remain in range.  The estimate
$6000(\log y)^3\le L^3$ holds eventually for every fixed cap
coefficient, because $\Xi^3=L^2\ell=o(L^3)$.

\subsection{Exact algebraic path-modulus bounds}
\label{rf:path}

Let $t=\rho^2=9/16$.  At depth $r$, with ancestors $p_j$ and
current prime $q$, the ladder gives uniformly
\[
 \log p_j\le(\tau t^j+o(1))\Xi\quad(0\le j<r),\qquad
 \log q\le(\tau t^r+o(1))\Xi.
\]
The cumulative special-box contribution is
$O(\Lambda\ell)=o(\Xi)$.  In the two-point shift of source
Lemma~4.4, put
$\Delta_j=|b_j^*a_1-a_j^*b_1|$.  The two moduli are
\begin{equation}
 m_1=3\prod_{j<r}\Delta_j
                 \frac{|b_2a_1-a_2b_1|}{q},\qquad m_2=qm_1.
 \label{rf:pathmoduli}
\end{equation}
Since $\Delta_j\le2E_{p_j}E_q$ and
$|b_2a_1-a_2b_1|/q\le E_q^2/q=\exp(O(\Lambda))$,
\[
 \log m_1\le(\tau F_r+o(1))\Xi,\qquad
 \log m_2\le(\tau G_r+o(1))\Xi,
\]
where
\[
 F_r=\frac12\left(\sum_{j<r}t^j+rt^r\right),\qquad
 G_r=F_r+t^r.
\]
Their exact differences are
\begin{equation}
 F_{r+1}-F_r=t^r\frac{25-7r}{32},\qquad
 G_{r+1}-G_r=t^r\frac{11-7r}{32}.
 \label{rf:fdiff}
\end{equation}
Thus
\begin{equation}
 \max F_r=F_4=\frac{40261}{32768},\qquad
 \max G_r=G_2=\frac{181}{128}.
 \label{rf:fmax}
\end{equation}

For an algebraic scale $X=Ma^d/q$ or $Ma^d$, tuning gives
$4u(\log y-\kappa_0)=\log X+O(1)$.  Since $q\ge B$,
\begin{align*}
 \log X&\ge(\mu+\delta\eta_0+o(1))\Xi+(d/2-1)\log q,\\
 \log Q_{\rm AP}&\ge\tfrac18\log X-\tfrac1{16}\log y+O(1)
             \ge(27/16+o(1))\Xi.
\end{align*}
For the last inequality, the coefficient of $\log q$ is eventually
positive, so minimizing at $q=B$ absorbs degree-rounding errors
into $O(\Lambda)$.  The exact worst-case comparison is
\begin{equation}
 \frac{27}{16}-\tau\max G_r
 =\frac{27}{16}-\frac{3439}{2048}
 =\frac{17}{2048}>0.
 \label{rf:algapgap}
\end{equation}
This replaces the source's displayed $0.51$ and $0.85$ envelopes;
their literal values do not hold for the new cap.

\subsection{A failed schedule and its fixed-root repair}
\label{rf:rootrepair}

There is a different progression in the root rational first moment
of source Lemma~4.3, p.~33.  Fix $(a,b)$ and $g_0$.
The allowed $g_1$ values form a progression of step $q_0$, and
\[
 z=(ag_1-bg_0)/q_0=z_{\min}+as.
\]
The modulus is $a$, rather than a pair determinant divided by
$q_0$.  Uniformly for $q_0\in(Q_0/2,Q_0]$,
\[
 \log a=(\tau/2+o(1))\Xi,\quad
 \log z=(g-\tau/2+o(1))\Xi,\quad
 \log Q_{\rm AP}=((g-\tau/2)/8+o(1))\Xi.
\]
Therefore this first moment requires $g>9\tau/2$.
Keeping $g=3$ with $\tau=19/16$ gives the negative margin
\begin{equation}
 \frac{3-\tau/2}{8}-\frac\tau2=-\frac{75}{256}.
 \label{rf:failedgap}
\end{equation}
That schedule is invalid even though \eqref{rf:algapgap} holds.
With $g=43/8$, the margin is instead
\begin{equation}
 \frac{g-\tau/2}{8}-\frac\tau2=\frac1{256}>0.
 \label{rf:rootapgap}
\end{equation}

The repair increases the root smoothness parameter:
\[
 \log z=(153/32+o(1))\Xi,\qquad
 \frac{4u}{U}=\frac{g-\tau/2}{18}+o(1)=\frac{17}{64}+o(1).
\]
Replace the source's root bound $4u\le0.16U$ by $4u\le0.27U$.
The lower admissibility margin is
$17/64-2/15=127/960>0$.
The density must now be estimated using its actual scale, not the
crude universal $\log(5u\log y)\le\ell+O(1)$:
\[
 u\log y=(153/128+o(1))\Xi,\qquad
 \log(5u\log y)=(2/3+o(1))\ell.
\]
Equation~\eqref{rf:densitylower} therefore gives
\begin{equation}
 \nu_{\rm rat}\ge
 \exp\{-(17/96+o(1))\Lambda\}.
 \label{rf:ratdensity}
\end{equation}
Since $E_q^2/q=\exp((4+o(1))\Lambda)$, the root first moment
satisfies
\begin{equation}
 \mu_{\rm rat}\ge
 \exp\{(4-17/96+o(1))\Lambda\}\ge e^{3.82\Lambda},
 \qquad 4-17/96-3.82=\frac7{2400}>0.
 \label{rf:ratmean}
\end{equation}
Thus the source's rationally good threshold
$|\mathcal R_q|\ge e^{3.80\Lambda}$ survives.

\subsection{Rounding, root fibres, and close pairs}

All root scales for fixed $q$ lie in
$[E_qG/(4q),2E_qG/q]$, an interval of logarithmic length $\log8$.
The ceiling defining $u$ has at most one jump in that range for
sufficiently large $N$.  On each constant-$u$ interval, the
unrounded grid index has derivative in $\log X$ of magnitude
\[
 \frac{10u_{\rm cap}}{4u}=\frac{2560}{17}+o(1).
\]
There are consequently fewer than $320$ ideal rounding
boundaries, preserving the source bound $1000$.
For the actual dyadic implementation, enlarge each boundary by
its logarithm-approximation error band.  No monotonicity of an
arbitrary approximation routine is assumed.  The error is at
most $\exp(-(60\log2+o(1))\ell^2)$, much smaller than the
$\exp(-30\ell^2)$ relative width of a $G$-block.  Outside these
bands the parameters agree with the ideal formulas.  Within a
single block, the possible $u$ values differ by at most one and
the grid indices for each fixed $u$ differ by at most one; the
at-most-four candidate parameter pairs are still valid.  This
band formulation is sufficient for the source's boundary-area
estimate.

Keep the root fibre cutoff
$T_{\min}=\lceil e^{-0.19\Lambda}L_G/E_q\rceil$,
where $L_G\asymp G/M_G$.
Using \eqref{rf:ratdensity}, its boundary contribution relative to
the square of the mean has coefficient at most
\[
 -0.19+17/96+o(1)=-31/2400+o(1)<-0.01.
\]
The same margin applies to the boundary term in the slope-bin
second moment.  On the long interior fibres, the original lower
interval bound is $H/X\ge e^{-0.195\Lambda}$; the factors $G$
and $E_q$ cancel.  On short-determinant fibres, the rough bound
uses $u'=2$, because $6D(z/H)\le9z/T_{\min}\le e^{4.21\Lambda}
\ll(y/2)^3$.  Thus no extension of the AP interval range
\eqref{rf:apconditions} is needed.  The source slope-bin conclusion
and its later root algebraic mean lower bound $e^{3.30\Lambda}$
retain strict margins.

For algebraic close pairs, increasing $M_{\rm pair}$ gives
frequency at most $e^{-0.18\Lambda}$.
Well-separated pairs still have $H/X\ge e^{-0.195\Lambda}$.
For the remaining close pairs choose
\begin{equation}
 u'=\left\lceil\frac{1.47\Xi}{\log(y/2)}\right\rceil+4.
 \label{rf:roughprefix}
\end{equation}
Here $1.47>\tau\max F_r$, and the extra four factors cover the
standalone $q$ and $B$ factors in the rough prefix conditions.
Indeed $\log q/\log y\le16/9+o(1)$ and
$\log B/\log y\le1/0.99+o(1)$, whose sum is below three.
Since $\log y\ge(17.82+o(1))\Lambda$,
$u'<0.083U<2u$ by \eqref{rf:algu-lower}.
Each rough weight loss is at most $e^{0.084\Lambda}$;
the two-field loss is at most $e^{0.17\Lambda}$.
After the polynomial comparability factors, the close-pair
contribution is at most $\exp((-0.01+o(1))\Lambda)$ times the
squared mean.  The
bad-child factor $e^{-0.93\Lambda}$ still leaves a bound stronger
than the source's $e^{-0.80\Lambda}$ target.

\subsection{Coefficient widths and sampling margins}

For the ancestor product together with two current evaluation
vectors, the coefficient on the $\Xi$ scale is at most
\[
 \tau\left(\frac12\sum_{j<r}t^j+t^r\right)
 =\tau\left(\frac87-\frac{t^r}{7}\right)
 <\frac{8\tau}{7}=\frac{19}{14}.
\]
This is below the fixed $34/25$ bound used in the coupling
obligation.  Including a possible modulus $q$ and the base scale,
the limiting tile-displacement coefficient is bounded by
\begin{equation}
 g+\tau+8\tau/7=\frac{887}{112}<8,
 \qquad 8-\frac{887}{112}=\frac9{112}>0.
 \label{rf:tilegap}
\end{equation}
With the fixed $34/25$ envelope the margin is still $31/400$.
Since $\Xi/\Lambda=U\longrightarrow\infty$, either positive
fixed margin makes the tile errors smaller than all required
$e^{-c\Lambda}$ errors.  Ancestor factors cancel in the ratio of
the inner and outer fibre boxes, so the tail argument's
$\Lambda$-scale combinatorial cost does not grow.

The global sampler satisfies
\[
 \log(W_{\rm cut}/G)=(13/8+o(1))\Xi,
 \qquad(d-J-4)\log G=(43/16+o(1))L>L.
\]
Its volume errors therefore remain smaller than the printed
$\RFLT{-0.98}$ error used before the total-variation bound.
The leading-coefficient and derivative shifts have width with
coefficient $8-g=21/8$, still much greater than the factor-base
scale.  Algebraic convolution shifts have width coefficient at
least $8-8\tau/7=93/14$, compared with a current coordinate of
coefficient at most $19/32$.  The extra $q$-preserving shifts
have width coefficient at least $8-\tau-8\tau/7=611/112$,
greater than $19/16$.  These replace, rather than reuse literally,
the source's displayed $3.01$, $4.20$, $4.9$, and $7.40$ envelopes.

The coefficient degree and height are unchanged.  Character primes
have size $x_{\rm ch}=2^{K_{\rm ch}}$ with
$K_{\rm ch}=\lceil e^{(8/15)\widehat\Lambda}\rceil$.
Thus $x_{\rm ch}\gg Q_0,M$ for the new cap as well.
The source's discriminant and character bit-size estimates need
no changed leading constant.  Their analytic certification remains
part of the contract.

\section{Sparse relation assembly and exact modular squares}
\label{rf:relations}

For $\iota=1,2$, let $K_\iota=\RFQ(\alpha_\iota)$,
$f_\iota(\alpha_\iota,1)=0$, and
$\theta_\iota=c_{\iota,d}\alpha_\iota$.
The polynomial
$\widetilde f_\iota(X)=c_{\iota,d}^{d-1}
 f_\iota(X/c_{\iota,d},1)$ is monic, and
$A_\iota=\RFZ[\theta_\iota]$ has the evaluation map
\[
 \phi_\iota:A_\iota\longrightarrow\RFZ/N\RFZ,
 \qquad\theta_\iota\longmapsto c_{\iota,d}m_N.
\]
Set $\gamma_\iota(a,b)=c_{\iota,d}a-b\theta_\iota$.
Then
\[
 \phi_\iota(\gamma_\iota(a,b))
 =c_{\iota,d}g_1^{-1}(ag_1-bg_0)\pmod N.
\]
The exact root identity and algebraic-node identity are
\begin{align*}
 q_0&\equiv c_{1,d}^{-1}g_1n_{\rm rat}^{-1}
                     \phi_1(\gamma_1(a,b))\pmod N,\\
 1&\equiv c_{\iota,d}c_{\iota',d}^{-1}
       \phi_\iota(\gamma_\iota(a,b))^{-1}
       \phi_{\iota'}(\gamma_{\iota'}(a,b))\pmod N.
\end{align*}
All inverted residues pass the unit checks; a failed inverse can
itself yield a proper gcd.

Propagating signs down the trees produces, for accepted trial $m$,
lists of primitive pairs with signs $\sigma_{m,\iota,\nu}\in\{\pm1\}$
and a relation
\begin{align}
 w_m&\equiv c_{1,d}^{e_{m,1}}c_{2,d}^{e_{m,2}}g_1^{e_{m,g}}
       \prod_{p\le B}p^{u_{m,p}}
       \phi_1(\Gamma_{m,1})\phi_2(\Gamma_{m,2})\pmod N,
 \label{rf:assembled}\\
 \Gamma_{m,\iota}&=\prod_\nu\gamma_\iota(a_{m,\iota,\nu},b_{m,\iota,\nu})
                            ^{\sigma_{m,\iota,\nu}}.
 \notag
\end{align}
The parity column has three scalar rows, rational-prime rows,
degree-one local-length rows in each field, and the sampled
quadratic-character rows.  Store only nonzero entries, labelled by
their type and prime or ideal data.  Sorting labels gives row
indices.  Signed exponents and the original factor lists are
retained to reconstruct \eqref{rf:assembled}; a parity column
alone is insufficient for square-root extraction.

Each relation contains at most
$2r_0\exp(O(\ell^2))$ nodes, each with polynomially many
prime-factor occurrences and character entries.  Therefore
\begin{equation}
 D_{\rm row},R_{\rm tot}=\exp((18+o(1))\Lambda),\qquad
 s_A=\exp((18+o(1))\Lambda),
 \label{rf:matrixsizes}
\end{equation}
where $s_A$ is the number of nonzero matrix entries.  This is an
actual sparse representation bound; no dense parity vector is
allocated and later discarded.

Given three independent binary null vectors, their combinations
map linearly into the joint obstruction space.  Under the
dimension-two event in \textup{(C3)}, some nonzero combination has
both $\Gamma_{T,\iota}=\prod_{m\in T}\Gamma_{m,\iota}$ square in
$K_\iota$.  Testing all seven combinations is a finite exact way
to find one, without computing the obstruction space itself.
For each field, put
\[
 P^-_{T,\iota}=\prod_{m\in T,\ \nu:\sigma_{m,\iota,\nu}=-1}
                         \gamma_\iota(a_{m,\iota,\nu},b_{m,\iota,\nu}).
\]
The positive product tested by the square-root routine is
\[
 \widetilde\Gamma_{T,\iota}
  =\widetilde f_\iota'(\theta_\iota)^2
      (P^-_{T,\iota})^2\Gamma_{T,\iota}
  =\widetilde f_\iota'(\theta_\iota)^2
       \prod_{m\in T,\nu}\gamma_\iota(a_{m,\iota,\nu},b_{m,\iota,\nu}).
\]
The derivative factor places a true square root in the polynomial
order as required by source Lemma~5.2.  Odd degree and irreducibility
modulo three permit initialization in $\RFF_{3^d}$ and quadratic
Newton--Hensel lifting.  The lifted polynomial is accepted only
after its exact square is checked in $\RFZ[X]/(\widetilde f_\iota)$.
Dividing the evaluated root by
$\widetilde f_\iota'(c_{\iota,d}m_N)\phi_\iota(P^-_{T,\iota})$
and multiplying by the half-exponents of the scalar and rational
parts in \eqref{rf:assembled} gives $Y$.
The final modular check $Y^2=\prod_{m\in T}w_m$ remains mandatory.

\section{A verified sparse rectangular kernel algorithm}
\label{rf:sparse}

The following result is independent of the factoring contract.
It is a specialization of classical black-box/Wiedemann linear
algebra \cite{Wiedemann1986}, with the three-vector deficiency and
deterministic space requirements made explicit.

\begin{theorem}[Sparse binary kernel with three outputs]
\label{rf:kerneltheorem}
Let $A\in\RFF_2^{a\times b}$, $b\ge a+3$, be stored as a list of
$s$ nonzero entries, and put $n=a+b$.  A finite randomized attempt
returns either failure or three verified independent vectors in
$\ker_{\RFF_2}A$.  It never returns incorrect vectors, succeeds
with probability at least $3/4$, and uses
\begin{align*}
 O\bigl(n(n+s)\operatorname{polylog}n\bigr)
 &\quad\text{bit operations},\\
 O\bigl((n+s)\operatorname{polylog}n\bigr)
 &\quad\text{bits of space}
\end{align*}
on a multitape machine.  A fixed number of attempts preserves
these bounds and gives the corresponding amplified success
probability.
\end{theorem}

\subsection{Preconditioning without a rank oracle}

If $A=0$, return three coordinate vectors.  Otherwise construct
$\RFF_Q$, where
\[
 h=\lceil\log_2(128n)\rceil,\qquad Q=2^h.
\]
An irreducible binary polynomial of degree $h$ can be found by
enumeration in $O(Q\operatorname{poly}(h))$ bit operations and
$\operatorname{poly}(h)$ space.  For example, test divisibility
of $X^{2^h}-X$ and the usual gcd conditions against
$X^{2^j}-X$; all powers are reduced modulo the candidate.
Since $Q=O(n)$ this initialization fits the theorem's bounds.

Choose independent uniform diagonal matrices
$D\in\RFF_Q^{b\times b}$ and $E\in\RFF_Q^{a\times a}$.
Zero diagonal values are allowed.  Define the implicit map
\begin{equation}
 P=ADA^{\mathsf T}E:\RFF_Q^a\longrightarrow\RFF_Q^a.
 \label{rf:preconditioner}
\end{equation}
It is applied using two sparse matrix-vector products and two
diagonal products, without forming its entries.

\begin{lemma}[Semisimple zero with high probability]
\label{rf:precondlemma}
Let $r=\operatorname{rank}A$.  With probability at least
$1-2r/Q$, the map $P$ has rank $r$, image $\operatorname{im}A$,
and
\[
 \RFF_Q^a=\ker P\oplus\operatorname{im}P.
\]
On this event its restriction to $\operatorname{im}A$ is invertible.
\end{lemma}
\begin{proof}
The coefficient of $z^r$ in $\det(I+zP)$ is
\[
 \sum_{|I|=|J|=r}
      \det(A_{I,J})^2\prod_{i\in I}E_i\prod_{j\in J}D_j.
\]
This follows by principal minors and Cauchy--Binet.  A rank-$r$
minor is nonzero, and different pairs $(I,J)$ give different
monomials, so the polynomial is nonzero even in characteristic
two.  Its degree is $2r$.  A nonzero polynomial of total degree
$k$ on independent uniform elements of $\RFF_Q$ vanishes with
probability at most $k/Q$; induction on the number of variables
proves this elementary polynomial zero bound.

On the nonvanishing event, rank $P$ is at least $r$, and it is
always at most rank $A=r$.  Its image is therefore
$\operatorname{im}A$.  The algebraic multiplicity of zero is
$a-r$, equal to its geometric multiplicity.  Thus zero has no
nontrivial Jordan block, giving the displayed direct sum and
invertibility on the image.
\end{proof}

The good event is not tested and the rank need not be computed.
It is used only to lower-bound the success probability.

\subsection{Krylov solution and exact verification}

Draw $v\in\RFF_Q^b$ uniformly and let $b_0=Av$.
If $b_0=0$, use $x=0$.
Otherwise draw a uniform row $w\in\RFF_Q^{1\times a}$ and compute
\[
 s_i=wP^ib_0\qquad(0\le i<2a).
\]
Berlekamp--Massey gives the monic minimal recurrence polynomial
$f(X)=\sum_{j=0}^e f_jX^j$ for this prefix, with $e\le a$.
If $f_0=0$, reject this attempt.  Otherwise form
\begin{equation}
 x=-f_0^{-1}\sum_{j=1}^e f_jP^{j-1}b_0,
 \label{rf:krylovsolve}
\end{equation}
and verify $Px=b_0$ exactly.  The powers are replayed in a fresh
pass; storing all Krylov vectors is unnecessary.

For an explicit recurrence implementation, start with connection
polynomials $C=B_{
m BM}=1$, order $e=0$, shift $m=1$, and
nonzero discrepancy $b_{
m BM}=1$.  At position $i$, compute
$\Delta=s_i+\sum_{j=1}^e C_js_{i-j}$.
If $\Delta=0$, increment $m$.
Otherwise save $T=C$ and replace
$C$ by $C-(\Delta/b_{
m BM})X^mB_{
m BM}$.
When $2e\le i$, also replace
$e$ by $i+1-e$, $B_{
m BM}$ by $T$,
$b_{
m BM}$ by $\Delta$, and $m$ by one;
otherwise increment $m$.
After $2a$ positions, return $f(X)=X^eC(X^{-1})$.
This is the standard finite-field recurrence algorithm; its
quadratic polynomial updates use $O(a^2)$ field operations and
$O(a)$ field elements of storage.

\begin{lemma}[Projection loses the true recurrence rarely]
\label{rf:krylovlemma}
Conditional on the good event in Lemma~\ref{rf:precondlemma},
for each fixed $b_0\in\operatorname{im}A$ the above solve succeeds
with probability at least $1-a/Q$.  When it succeeds,
\eqref{rf:krylovsolve} is the unique solution in
$\operatorname{im}A$.
\end{lemma}
\begin{proof}
For $b_0\ne0$, its true Krylov minimal polynomial $f_*$ has degree
at most $a$ and has nonzero constant term, because $P$ is
invertible on its image.  If a scalar projection loses a factor
of $f_*$, then for some distinct irreducible factor $g$ of $f_*$,
the scalar sequence is annihilated by $f_*/g$.  In particular,
\[
 w\,(f_*/g)(P)b_0=0.
\]
The vector after $w$ is nonzero by minimality of $f_*$.
Each such event has probability $1/Q$, and there are at most $a$
distinct irreducible factors.  Their union has probability at
most $a/Q$.

A prefix of length $2a$ suffices: a candidate recurrence of order
at most $a$ which annihilates that prefix gives a difference
sequence with at least $a$ initial zeros.  That difference
sequence obeys the true order-at-most-$a$ recurrence, hence is
identically zero.  On the no-loss event the returned polynomial
is $f_*$.  The identity $f_*(P)b_0=0$ proves
$Px=b_0$.  The polynomial expression for $x$ lies in
$\operatorname{im}A$, where the solution is unique.  Even in the
loss event, the explicit verification prevents an incorrect
linear solve from being used.
\end{proof}

Set
\begin{equation}
 k=v-DA^{\mathsf T}Ex
 \label{rf:kernelprojection}
\end{equation}
and verify $Ak=0$.
On the good preconditioner event, the exact-image solution makes
\eqref{rf:kernelprojection} a linear projection onto $\ker A$;
it is the identity on that kernel.  Thus its ideal output for
uniform $v$ is uniform on $\ker_{
 \RFF_Q}A$.

Choose a fixed binary basis $e_1,\ldots,e_h$ of $\RFF_Q$ and
write $k=\sum_{j=1}^h e_jk_j$.
Since $A$ is binary, every $k_j$ lies in
$\ker_{
 \RFF_2}A$.  A uniform element of the extension-field kernel has
independent uniform binary components in that kernel.  Its
dimension is at least three.  Projecting onto any three-dimensional
quotient and taking a union bound over the seven nonzero binary
linear forms shows
\[
 \RFPr\{\operatorname{span}(k_1,\ldots,k_h)
                  \text{ has dimension less than three}\}
 \le7\,2^{-h}=7/Q.
\]
Extract three independent components by binary elimination and
verify each matrix product and their independence.

\begin{proof}[Proof of Theorem~\ref{rf:kerneltheorem}]
Compare the algorithm with the ideal exact projection for each
good preconditioner.  Its failure probability is at most
\[
 \frac{2r}{Q}+\frac aQ+\frac7Q
 \le\frac{3a+7}{Q}<\frac14.
\]
This comparison is made before conditioning on verified success;
no claim of uniformity after a biased acceptance event is needed.
All returned vectors are checked, so there is no false success.

There are $O(a)$ applications of $P$ for the sequence and the
polynomial replay.  Field operations use
$\operatorname{poly}(h)$ bit work.  To multiply a sparse matrix
by a vector on a multitape machine, sort its index/value records,
merge them with the vector entries, then sort and sum by output
index.  Sequential merge sort uses
$O((n+s)\operatorname{polylog}n)$ bit work and space per product.
The vector, diagonal, recurrence, and component data have at
most $O(n\operatorname{polylog}n)$ bits; no dense $P$ is formed.
The recurrence's $O(a^2)$ field operations fit the time bound.
Every loop is finite, proving the deterministic space bound.
\end{proof}

Applying the theorem to \eqref{rf:matrixsizes} costs
$\exp((36+o(1))\Lambda)$ time and
$\exp((18+o(1))\Lambda)$ space per attempt.
Three independent attempts fail with probability at most $1/64$.
This factor is harmless in the phase success probability.

\section{Hidden-root symmetry and streaming}
\label{rf:blindness}

The permission to discard rejected roots and use randomized
linear algebra follows from an exact finite-group fact.

\begin{theorem}[Several blind dependencies]
\label{rf:blindtheorem}
Let $G$ be a finite abelian group, $h:G\to G$ a homomorphism,
and $H=\ker h$.  Draw $z_1,\ldots,z_m$ independently uniformly
from $G$.  Reveal $h(z_i)$ and an independent random tape.
Conditional on that transcript, let an integer matrix
$E\in\RFZ^{r\times m}$ and elements $Y_1,\ldots,Y_r$ be fixed,
with
\[
 h(Y_j)=\prod_i h(z_i)^{E_{ji}}.
\]
Then $(\prod_i z_i^{E_{ji}}Y_j^{-1})_{j=1}^r$ is uniform on a
coset of the image of the homomorphism $E:H^m\to H^r$.
For squaring, if $H$ is elementary abelian two-torsion and
$E\bmod2$ has full row rank, this image is $H^r$.
\end{theorem}
\begin{proof}
Conditional on the transcript, each $z_i$ is uniform on its
$H$-coset, independently of the others.  A homomorphism of finite
groups maps the uniform distribution onto the uniform distribution
on its image, because every nonempty fibre is a coset of its
kernel.  If $H\cong\RFF_2^k$, a full-row-rank binary matrix induces
a surjection on each of the $k$ coordinates, proving the last
assertion.
\end{proof}

For odd $N$ with $k\ge2$ distinct prime divisors,
$H=\{\varepsilon\in(\RFZ/N\RFZ)^*:\varepsilon^2=1\}$
has $2^k$ elements, including when $N$ has repeated prime factors.
For a single verified nonempty dependency, $ZY^{-1}$ is uniform
on $H$.  The gcd is trivial exactly at the two global signs,
so its split probability is $1-2^{1-k}\ge1/2$.
For $r$ independent verified dependencies selected before any
unsealing, the conditional probability that all give trivial
gcds is $(2^{1-k})^r$.

The main algorithm needs only one dependency.  Keeping three
independent null vectors and checking their seven combinations
is for the dimension-two square-class obstruction, not for an
unjustified independence assumption about chosen relations.
If several successful dependencies are desired, collecting
$s+r$ extra columns when the obstruction dimension is at most
$s$ yields a good-combination subspace of dimension at least $r$.
For fixed $s+r$, exhaustive exact tests add only a constant factor.
Every test and selection must finish before any root is unsealed.

Only the current trial's root and roots of accepted relations need
be stored.  Acceptance, sparse preconditioning, Krylov projections,
and square-root selection depend only on exposed squares and
independent coins.  Erasing rejected roots therefore does not
change the conditional coset law.  A later failure must not be
used to adapt a supposedly blind selection using roots that have
already been revealed; a fresh phase uses fresh roots.

\section{Collection probability and the exact time recurrence}
\label{rf:recurrence}

\subsection{Fewer, larger root primes}

The source proves a parameter-general rational smoothing
proposition, Proposition~2.3(2).  In the rough-input regime used
here, a shorter counting argument gives the needed estimate
independently.

\begin{lemma}[Square-class counting for rough inputs]
\label{rf:rough-ratios}
Let $N$ be odd, let $\mathcal Q$ be a set of $m$ distinct primes
coprime to $N$ and at most $Q_0$, and let $r\ge1$ satisfy
$Q_0^r\le\sqrt N/2$.  Put
$H=|(\RFZ/N\RFZ)^*/Q_N|=2^{\omega(N)}$.
If $Hr^2/m\le1/2$, then a uniform $w\in Q_N$ has probability
at least
\[
 \frac1{2\varphi(N)}\binom mr^2
\]
of having its normalized shortest lattice vector equal to the
products of two disjoint $r$-element subsets of $\mathcal Q$.
For the $B$-rough inputs and good-prime sets of this manuscript,
the condition $Hr_0^2/|\mathcal Q_{\rm good}|\le1/2$ holds
eventually.
\end{lemma}
\begin{proof}
Partition the $r$-element subsets by the square class of their
prime product.  If the class sizes are $c_1,\ldots,c_H$, then
the number of ordered pairs in the same class is
\[
 \sum_jc_j^2\ge H^{-1}\left(\sum_jc_j\right)^2
                  =H^{-1}\binom mr^2.
\]
For each first subset, the number of second subsets intersecting
it is at most $r\binom{m-1}{r-1}$.
The total number of intersecting ordered pairs is therefore at
most $(r^2/m)\binom mr^2$.  Subtracting all of them leaves at
least $(2H)^{-1}\binom mr^2$ disjoint pairs in the same class.

Write their products as $x,y$.  They are coprime units, and
$x/y\bmod N$ is a square.  Different pairs give different ratios:
congruence of two ratios makes $xy'-x'y$ a multiple of $N$ of
absolute value less than $N$, hence zero; coprimality and unique
factorization recover the same two subsets.  For such a ratio,
the vector $(x,y)$ has squared norm at most $N/2$.  A shortest
vector $(x_1,y_1)$ in the same lattice has determinant
$x_1y-y_1x$ divisible by $N$ and absolute value at most $N/2$,
so it is parallel to $(x,y)$.  Primitivity and shortestness make
it equal to $(x,y)$ up to sign; the normalization $x_1>0$
chooses the positive pair.  Dividing the pair count by $|Q_N|$
and using $H|Q_N|=\varphi(N)$ proves the bound.

For a $B$-rough input,
$\log H\le(\log2)L/\log B=o(\Xi)$.
The good-prime set has logarithmic cardinality
$(\tau+o(1))\Xi$, while $\log r_0=O(\ell)$.
Thus $Hr_0^2/|\mathcal Q_{\rm good}|\longrightarrow0$.
\end{proof}

The prime number theorem and \textup{(C1)} give
$m=|\mathcal Q_{\rm good}|\ge3Q_0/(10\log Q_0)$ eventually.
Using $\binom mr\ge(m/r)^r$, $\varphi(N)\le N$, and
$Q_0^{2r_0}\ge N/8$, the lemma gives
\[
 p_{\rm ratio}\ge\frac1{16}
     \left(\frac{10r_0\log Q_0}{3}\right)^{-2r_0}.
\]
Since $r_0\longrightarrow\infty$, this implies the convenient
weaker bound
\begin{equation}
 p_{\rm ratio}\ge
 \frac18(4r_0\log Q_0)^{-2r_0}.
 \label{rf:ratiobound}
\end{equation}
Excluding divisors of the fixed setup integers removes only
polynomially many primes.  The bound refers
to the actual shortest-vector rational-smoothing interface, not
to independent uniform numerator and denominator samples.

Here
\[
 2r_0=(16/19)U+O(1),\qquad
 4r_0\log Q_0=2L+O(1),
\]
so \eqref{rf:ratiobound} gives
\begin{equation}
 p_{\rm ratio}\ge
 \exp\{-(16/19)\Lambda-O(U+\ell)\}.
 \label{rf:acceptance}
\end{equation}
The union bound over $2r_0$ independent fresh
tree attempts loses only a factor tending to one, since each
eligible prime fails with probability at most $L^{-2}$.

The $\epsilon_N\Lambda=U\sqrt\ell$ term in
\eqref{rf:trialcap} dominates $O(U+\ell)$.
Conditional on a good setup, independent capped trials therefore
have mean accepted count at least $16R_{\rm tot}$ eventually.
Chebyshev's inequality, or a binomial tail bound, makes the
probability of collecting fewer than $R_{\rm tot}$ at most
$1/100$ for all sufficiently large $N$.
If a nonunit draw occurs before enough relations are collected,
its gcd already splits $N$.  Conditioning on the unit-indicator
vector leaves independent uniform units before that first event,
exactly as required by Theorem~\ref{rf:blindtheorem}.

The logarithmic number of expensive root trials is consequently
\begin{equation}
 \log M_{\rm cap}=(A+o(1))\Lambda,
 \qquad A=18+\frac{16}{19}=\frac{358}{19}.
 \label{rf:A}
\end{equation}
For comparison, the old $\tau=1/3$ gives coefficient
$18+3=21$, after removing the source's fixed cap slack.

\subsection{A phase has constant success probability}

Under the contract, multiply the setup probability $3/4$, the
character probability $98/100$, collection probability $99/100$,
and three-attempt sparse-kernel probability $63/64$.
The dimension-two argument provides a certified square
dependency, and the sealed-root theorem splits with conditional
probability at least $1/2$.  The product exceeds $1/3$; all
remaining capped-draw errors tend to zero and can be absorbed
while retaining a positive absolute lower bound.  Using a slightly
larger effective threshold retains the source's at-most-three
expected phases if desired.  Only a fixed positive phase
probability is needed for the leading-constant analysis.

\subsection{Solving the recursive costs}

Suppose the strong-induction majorant for complete factorization
has coefficient $C$.  Each root trial makes at most two recursive
calls on integers at most $\sqrt N/2$.  Since
\[
 \frac{\Lambda(N^\alpha)}{\Lambda(N)}
 =\alpha^{1/3}
   \left(1+\frac{\log\alpha}{\ell}\right)^{2/3}
 =\alpha^{1/3}+o(1)
\]
for fixed $0<\alpha<1$, the root-call cost has coefficient
\begin{equation}
 A+2^{-1/3}C.
 \label{rf:rootrec}
\end{equation}
There are at most $2r_0M_{\rm cap}$ descent trees and
$\exp(O(\ell^2))$ nodes per tree.  The node trial caps are at
most $\exp(2.05\Lambda)$, and each trial makes only polynomially
many recursive calls.  Proposition~\ref{rf:treebounds} therefore
gives the second coefficient
\begin{equation}
 A+2.05+(19/64)^{1/3}C.
 \label{rf:descrec}
\end{equation}
The determinant and rational calls are smaller; the displayed
bound safely includes them.

Factor-base enumeration costs at most
$\exp((36+o(1))\Lambda)$, even when roots modulo each prime are
found by exhaustive evaluation.  Sparse linear algebra costs
$\exp((36+o(1))\Lambda)$, and the exact square-root extraction
bound is below coefficient $38$.  Character preparation and
all polynomial overhead fit below that same threshold.
It suffices that
\begin{equation}
 C>\max\left\{
 38,\quad \frac{358/19}{1-2^{-1/3}},\quad
 \frac{358/19+41/20}{1-(19/64)^{1/3}}
 \right\}.
 \label{rf:threshold}
\end{equation}
The two recursive thresholds are approximately $91.33375329$
and $62.75798$.  Take the rational value
$C=457/5=91.4$.
The exact certificates are
\begin{align}
 2(C-A)^3-C^3&=\frac{370059927}{857375}>0,
 \label{rf:cert1}\\
 64(C-A-41/20)^3-19C^3
             &=\frac{6795396287504}{857375}>0.
 \label{rf:cert2}
\end{align}
All quantities cubed are positive, so these inequalities are
equivalent to the required strict recurrence inequalities.
Numerically their left-hand cost coefficients before comparison
with $C$ are about $91.38633334$ and $81.86508294$.

\begin{proof}[Time part of Theorem~\ref{rf:main}]
Use strong induction on the integer input.  Exact preprocessing
and proper splitting imply only $O(\log N)$ unresolved worklist
items need be split.  A phase on any such item has the three
cost bounds above, uniformly over its tuple.  Every recursively
factored integer is smaller, and the induction applies to its
complete factorization, not just to one split.

The strict gaps in \eqref{rf:threshold} absorb all polynomial
factors, $\exp(O(\ell^2))$ factors, the expected constant number
of phases, step-for-step background interleaving, and the
uniform $o(1)$ terms.  Explicitly, choose a sufficiently large
effective threshold so their total coefficient is below $C$ by
a fixed positive amount.  A finite multiplicative constant
handles all smaller inputs, which terminate by trial division.
The induction majorant may be taken as
$K\operatorname{poly}(\log N)\exp(C\Lambda(N))$.
This proves the asserted logarithmic limsup bound.
\end{proof}

The recurrence allows every fixed $C$ above its dominant threshold;
$91.4$ is the convenient rational value certified in
\eqref{rf:cert1}--\eqref{rf:cert2}.  This calculation is a
conditional upper-bound analysis, not a lower bound for the source
algorithm or for factoring.

\section{Deterministic workspace}
\label{rf:space}

Every accepted relation has a factor list and sparse column of
size $\exp(o(\Lambda))$.  Equation~\eqref{rf:matrixsizes} therefore
bounds the whole relation collection, its sealed accepted roots,
and its sparse parity data by $\exp((18+o(1))\Lambda)$ bits.
An accepted root has only $O(L)$ bits; a rejected root is erased.
The trial counter has $O(\Lambda)$ bits even though its numerical
cap is exponentially large on that scale.

The factor base occupies $O(dB\log N)=\exp((18+o(1))\Lambda)$
bits.  Character data and their arithmetic use the smaller bound
in \textup{(C3)}.  The sparse kernel theorem uses
$\exp((18+o(1))\Lambda)$ bits.  Across the fixed number of
tested dependencies, the positive factor-list length satisfies
\[
 M_{\rm fac}\le R_{\rm tot}\,2r_0\exp(O(\ell^2))
              =\exp((18+o(1))\Lambda),
\]
so \textup{(C4)} also uses
$\exp((18+o(1))\Lambda)$ bits.  Sequential multiplication and
replay are used rather than retaining every intermediate product.

Each recursive call operates on an integer at most $\sqrt N/2$
by an exact guard.  Recursive calls are executed sequentially,
reusing the child's workspace when it returns.  The active
stack has depth $O(\log\log N)$, and the sum of its per-level
spaces is $\exp((18+o(1))\Lambda(N))$.
For a formal slack argument, bound each sufficiently large
phase by $e^{18.5\Lambda}$ and compare it with the majorant
$e^{19\Lambda}$; the first child has coefficient at most
$19\,2^{-1/3}<19$.  Strong induction absorbs the finite small
range and all stack overhead.
Every random loop has a finite cap and failed phases overwrite
their storage.  Thus the space bound is deterministic over coin
tapes, completing the space part of Theorem~\ref{rf:main}.

\section{Limits of two proposed shortcuts}
\label{rf:shortcuts}

The following elementary screens explain why two inexpensive
modifications do not automatically remove the top recursive
factoring stage.  They concern the present sampler, not arbitrary
factoring algorithms.

\begin{proposition}[Asymmetric rational boxes]
For a uniform square $w$ in the square subgroup $Q_N$ of
$(\RFZ/N\RFZ)^*$, the probability that some unit ratio $x/y=w$
has $1\le x\le X$ and $1\le y\le Y$ is at most $XY/|Q_N|$.
For a $B$-rough $N$, a success probability
$\exp(-O(\Lambda))$ therefore requires
$XY\ge N^{1-o(1)}$; at least one coordinate bound is
$N^{1/2-o(1)}$.
\end{proposition}
\begin{proof}
There are at most $XY$ possible ordered pairs and hence at most
that many ratios.  For odd $N$,
$|Q_N|=\varphi(N)/2^{\omega(N)}$.
Roughness gives $\omega(N)\le L/\log B$ and
$N/\varphi(N)=1+o(1)$, so $|Q_N|=N^{1-o(1)}$.
\end{proof}

This does not exclude a partial factorization routine, a different
sampling law, or an algorithm avoiding complete factorization
of both coordinates.  It excludes the argument that merely
skewing the current box makes both recursive inputs much smaller.

\begin{proposition}[Accidental reuse of eligible root primes]
For $T=\exp(O(\Lambda))$ independent blind trials, each retaining
at most two positive coordinates bounded by $\sqrt N/2$, the
probability that an eligible prime $q\in(Q_0/2,Q_0]$ divides
coordinates in two different trials is
\[
 \exp\{-(\tau+o(1))\Xi\}=o(1).
\]
\end{proposition}
\begin{proof}
For a fixed $q$, there are $O(N/q)$ possible coordinate pairs
with $q$ dividing either coordinate.  Uniformity of the exposed
square bounds the probability in one trial by
$O(N/(q|Q_N|))$.  Independence and a union bound over trial
pairs and primes, using
$\sum_{q>Q_0/2}q^{-2}=O(1/Q_0)$ even over all integers, give
\[
 O\left(\frac{T^2}{Q_0}
            \left(\frac{N}{|Q_N|}\right)^2\right).
\]
Here $\log(N/|Q_N|)=O(L/\log B)=o(\Xi)$ and
$\log T=O(\Lambda)=o(\Xi)$, proving the claim.
Filtering to accepted trials cannot create a collision absent
from the complete presampled list.
\end{proof}

Thus a product-tree batch gcd of these same blind candidates
almost never obtains an eligible root prime by accidental
reuse.  This leaves open deliberately correlated samplers and
batching of different descent cofactors.
Likewise, generating already-factored numerator and denominator
products does not supply a square root with the required hidden
CRT-sign distribution.  That additional construction cannot be
assumed free.

\section{Why the first parameter is still one third}

Changing fixed constants $g$ and $\tau$ cannot change the first
$L$-parameter.  There is also a useful restricted balance
calculation.  Suppose one retains this explicit factor-base and
polynomial-value architecture, writes the degree as
$L^{a+o(1)}$ and $\log B=L^{b+o(1)}$, and charges reciprocal
smoothness density by the same method.  Three charged exponents
remain:
\begin{enumerate}[leftmargin=*]
\item the explicit factor base contributes $b$;
\item bottom-level polynomial norms have a smoothness parameter
of order the degree, contributing $a$;
\item the modular coefficient construction requires
$d\log G\gtrsim L$ and $\log M\ge\log G$.
Smoothing this coefficient height over $B$ contributes
$1-a-b$ when that quantity is positive.
\end{enumerate}
Their maximum obeys
\[
 \max\{a,b,1-a-b\}\ge\frac13.
\]
This is a limitation of that charging argument, not a complexity
lower bound for factoring.  It does not rule out a new relation
constructor, an implicit factor base with genuinely cheaper
access, a different coefficient-height mechanism, or a different
smoothness argument.  None of those missing constructions is
supplied by the enlarged-cap schedule.

\section{Reproducibility and unresolved obligations}
\label{rf:reproducibility}

The repository contains the exact rational verifier
\texttt{verify\_recursion\_budget.py}, the pinned-source metadata,
and the audit notes from which the present manuscript was
assembled.  Run the verifier with the repository's documented
Python command.  It uses only integer and rational arithmetic;
there is no floating-point decision in a decisive comparison.
It checks:
\begin{enumerate}[leftmargin=*]
\item the path difference identities and exact maxima in
\eqref{rf:fdiff}--\eqref{rf:fmax};
\item the algebraic conductor margin $17/2048$;
\item the failed $g=3$ root margin $-75/256$ and repaired margin
$1/256$;
\item the rational-root admissibility, density $17/96$, mean
margin $7/2400$, and boundary margin beyond $0.01$;
\item the fibre-width margin $9/112$, global sampler margin
$13/8$, and derivative-shift width $21/8$;
\item the rough-prefix and recurrence certificates, including
\eqref{rf:cert1}--\eqref{rf:cert2}.
\end{enumerate}
Finite checks of the path identities are accompanied by their
symbolic difference proof above, which establishes the maxima
for every depth.  The script does not prove an AP asymptotic,
an infinite probability assertion, or a complete factorer's
expected running time.  A verified factor output would likewise
not establish the expected-time analysis.

The main remaining obligations are precisely the enlarged
analytic and coupling statements in
Section~\ref{rf:analyticcontract}, their conditional propagation
through the entire adaptive descent, and an independent audit
of the imported algebraic certification.  An implementation must
also replace all hard-coded proof envelopes identified above;
editing only $\tau$ in a source program would be insufficient.
No full end-to-end implementation or formal proof-assistant
certificate is claimed by this manuscript.

The independent results supplied here are the sparse kernel
theorem, the hidden-root theorem, the rough-input square-class
counting lemma, the explicit finite parameter
comparisons, and the conditional recurrence and storage
deductions.  Their role is to turn a promising modification
into a precise, testable mathematical proposal while keeping
the unresolved proof contract visible.
